\documentclass[12pt]{article}
\usepackage[margin=1in]{geometry}
\usepackage{setspace}
\usepackage{titlesec}
\usepackage{enumitem}
\usepackage{hyperref}

\titleformat{\section}{\large\bfseries}{\thesection}{1em}{}
\titleformat{\subsection}{\normalsize\bfseries}{\thesubsection}{1em}{}

\hypersetup{
  colorlinks=true,
  linkcolor=blue,
  urlcolor=blue,
  pdftitle={Machine Vision Project Proposal},
  pdfauthor={Graduate Student Team},
  pdfsubject={Machine Vision Course Proposal},
  pdfkeywords={Machine Vision, Computer Vision, Image Processing, Project Proposal}
}

\begin{document}

\begin{center}
  {\LARGE \textbf{Machine Vision Project Proposal}}\\[1em]
  {\large \textbf{Course:} Machine Vision (Graduate Level)}\\
  {\large \textbf{Instructor:} Dr. Yuqian Zhang}\\
  {\large \textbf{Semester:} Fall 2025}\\[2em]
\end{center}

\noindent
\textbf{Team:} \underline{\hspace{12cm}}\\[1em]
\textbf{Project Title:} \underline{\hspace{12cm}}\\[1em]
\textbf{Keywords:} \underline{\hspace{12cm}}\\[2em]

\section*{Proposal Description}

\subsection*{1. Problem Statement}
\vspace{0.5em}
\noindent
This project addresses the challenge of performing \textbf{3D facial modeling and biometric verification using only standard RGB cameras}, without the need for specialized depth or infrared hardware. Current systems often rely on expensive and bulky sensors to capture detailed 3D information, which limits their accessibility and practical deployment.\\

This limitation motivates the development of novel algorithms and techniques that can infer accurate 3D facial structures and verify identities solely from conventional 2D images. Achieving this goal involves overcoming difficulties such as varying lighting conditions, pose variations, and the inherently ambiguous nature of monocular depth estimation.\\

The core problem is to design a robust and efficient method that enables reliable 3D facial recognition and verification using ubiquitous RGB imaging devices, thereby expanding the applicability of biometric technologies to a wider range of environments and users without specialized sensing.\\

\vspace{2em}

\subsection*{2. Background}
\vspace{0.5em}
\noindent
Provide relevant background information, including a brief review of existing methods, technologies, or research related to the problem. Highlight gaps or limitations that justify the need for your project.

\vspace{2em}

\subsection*{3. Objective}
\vspace{0.5em}
\noindent
State the primary goals and objectives of your project. What do you intend to achieve by the end of this work?

\vspace{2em}

\subsection*{4. Approach}
\vspace{0.5em}
\noindent
Outline the methodology and techniques you plan to use to tackle the problem. Include details about data sources, algorithms, tools, and any experimental design considerations.

\vspace{2em}

\subsection*{5. Expected Outcome}
\vspace{0.5em}
\noindent
Describe the anticipated results or deliverables of the project. Explain how these outcomes will demonstrate the success of your approach.

\vspace{2em}

\subsection*{6. Significance}
\vspace{0.5em}
\noindent
Discuss the importance and potential impact of your project within the field of machine vision. Include possible applications, benefits, and contributions to research or industry.

\vspace{3em}

\noindent
\textbf{Prepared by:} \underline{\hspace{6cm}} \hfill \textbf{Date:} \underline{\hspace{4cm}}

\end{document}
